\subsection{One top set for all large-order orbits}\label{orb:section}

Fix \(0<\varepsilon<1/52\) and a positive integer \(o\). Put
\[
 h=2+8\varepsilon,\qquad \gamma=\varepsilon/8,\qquad p=x^{-\gamma}.
\]
Let \(H\) be a finite abelian group of order \(n\asymp x^h\), with a
subgroup \(A_0\) of index \(o\). Fix \(T_2\subset A_0\) before choosing
the tops, and assume \(n_2=|T_2|\ge c n x^{-\varepsilon}\) for a
fixed \(c>0\). A thick set satisfying the reciprocal subpower bound of
Proposition~\ref{host:behrend} also satisfies this inequality after division
of its cardinality by any fixed positive constant, for all sufficiently
large \(x\).

The original orbit predicate uses the integers
\[
\begin{gathered}
 K=\lceil(h+2)/\gamma\rceil,\quad r=K+1,\quad k_0=6r+1,\quad
 U=\lceil1/p\rceil,\\
 \ell=k_0 U,\quad M=\lceil (h+3)10^{k_0}\log x\rceil,\quad
 G'=(M+5)\ell,\\
 N_1=\left\lceil\max\{6\exp(1)G'p,(h+3)\log_2 x\}\right\rceil+1 .
\end{gathered}
\]
Let \(L_2\) be a nonnegative integer with \(L_2=O_\varepsilon(\log x)\).
On a page of rank \(k\), one may take any upper bound for the original
\(L_2=2\sup_q\lceil\log_q((k x^{26})^4)\rceil\), where \(q\) runs
over the inert primes. Indeed \(q\ge2\) and \(k\le2x^{9/2+10\varepsilon}\).

For \(g\in H\) of order \(d\) and an origin \(a\in H\), label
\(a+ig\) by \(i=0,\ldots,d-1\). A cyclic window of length \(L\) starts
at \(i\) and consists of \(i,i+1,\ldots,i+L-1\bmod d\). The complete
tiles are the consecutive intervals of length \(\ell\) in
\([0,\lfloor d/\ell\rfloor\ell)\). In each tile number its \(k_0\)
consecutive cells of length \(U\) starting at one. A tile is marked
when its odd-numbered cells are empty and each even-numbered cell
contains exactly one top, at an even coordinate \(i\).

\begin{proposition}[Simultaneous orbit and overlap selection]\label{orb:selection}
For all sufficiently large \(x\) there is one set \(T\subset H\)
such that, for every \(g\) of order \(d\ge x^{2\gamma}\) and every
origin \(a\), all three assertions hold.
\begin{enumerate}
\item Every cyclic \(L_2\)-window contains at most \(K\) tops.
\item Starting at any complete tile, one of the next \(M+1\) complete
tiles, counting the starting tile, is marked. Here the tile indices
are taken modulo \(\lfloor d/\ell\rfloor\).
\item Every cyclic \(G'\)-window contains at most \(N_1\) tops.
\end{enumerate}
Write \(n_T=|T|\), let \(p_z=p\) on \(A_0\) and \(p_z=3p/4\) on
the other cosets, and put \(\lambda_C=(3+\mathbf{1}_{C=A_0})/(3o+1)\).
The same choice satisfies
\begin{align}
  pn/2\le n_T\le 2pn,\qquad
  \big||T\cap C|-\lambda_C n_T\big|
      &\le n_T x^{-1+\varepsilon} \quad(C\in H/A_0), \label{orb:cosets}\\
 \big||T\cap(z+T_2)|-p_z n_2\big|
      &\le (pn_2)^{1/2} x^{4\varepsilon}\quad(z\in H).
       \label{orb:overlap-count}
\end{align}
The threshold may depend on the constants in the size bounds for
\(n,n_2,L_2\), on \(o\), and on \(\varepsilon\).
\end{proposition}

\begin{proof}
Select the points of \(H\) independently, with inclusion probabilities
\(p_z\). For each qualifying direction \(g\) and each \(z\in H\),
make an event for excess occupancy in the \(L_2\)-window starting at
\(z\), and another for excess occupancy in the \(G'\)-window starting
there. For each of the two possible parities make one more event as
follows. Read \(\lfloor M/2\rfloor\) consecutive length-\(\ell\) tiles
starting at \(z\) in direction \(g\), in the physical group. Use the
specified parity for the acceptable local offsets from \(z\). The event
is that every tile in this string fails to be marked. There are at most
\(4n^2\) events. Each uses at most \(G'=O_\varepsilon(x^\gamma\log x)\)
different coordinates. Indeed \(L_2\le G'\) and \(G'<x^{2\gamma}\)
eventually. In particular every string used in the probability estimates
has distinct physical coordinates.

The short-window event has probability at most
\(\binom{L_2}{K+1}p^{K+1}\) by a union over subsets. Since
\((K+1)\gamma\ge h+2+\gamma\), this is at most
\(q_*=x^{-h-2-\gamma/2}\) eventually. For the long window the same
argument bounds the probability by
\(\binom{G'}{N_1}p^{N_1}\le(\exp(1)G'p/N_1)^{N_1}\le6^{-N_1}\)
when \(N_1\le G'\). The event is empty when \(N_1>G'\). The definition
of \(N_1\) gives \(6^{-N_1}\le x^{-(h+3)\log_2 6}\le q_*\).

For an odd-numbered cell the empty probability is at least
\((1-p)^U\), which tends to \(\exp(-1)\). For an even-numbered cell
at least \(\lfloor U/2\rfloor\) offsets have the prescribed parity.
The events that one particular such point is its unique top are
disjoint. Their total probability is at least
\[
 \lfloor U/2\rfloor(3p/4)(1-p)^{U-1}\longrightarrow 3/(8\exp(1)).
\]
Both limits exceed \(1/8\). Disjoint cells are independent, so each
tile is marked with probability at least \(8^{-k_0}\). Disjoint tiles
are independent as well. A failure string has probability at most
\[
 \exp(-8^{-k_0}\lfloor M/2\rfloor)
 \le \exp(1)x^{-(h+3)(10/8)^{k_0}/2}\le q_* .
\]
Here \(k_0\ge7\), \(h=2+8\varepsilon\), and \(0<\varepsilon<1/52\).
These bounds hold for both parity choices and every starting point.

Two events are joined if their sets of coordinates meet. A fixed
coordinate can occur in at most \(4G'\) events per direction, since
each possible position in the string determines its start. The
dependency degree is therefore at most \(D=4n(G')^2\), and
\(Dq_*=O_\varepsilon(x^{-2+3\gamma/2}(\log x)^2)=o(1)\).
For completeness, put \(q=2q_*\), so that
\(q_*/(1-q)^D\le q\) for large \(x\). The probability of any event
conditioned on avoiding any collection of the others is at most \(q\).
Induct on the size of the collection. Separate its non-neighbours
from its neighbours. The event is independent of the joint avoidance
of the non-neighbours by disjointness of the coordinates. When at least
one neighbour is present, the additional conditional probability of
avoiding the neighbours is at least \((1-q)^D\) by the induction
hypothesis and the chain rule. If no neighbour is present, independence
gives the claim directly. Multiplying the successive avoidance
probabilities now gives
\[
 \Pr(\text{all local events avoided})
 \ge (1-2q_*)^{4n^2}
 \ge \exp\big(-C x^{8\varepsilon-\gamma/2}\big).
\]
The last step uses \(n^2q_*\ll x^{h-2-\gamma/2}\).

Avoidance implies the three orbit assertions also at a tile-index wrap.
There are more than \(M+1\) complete tiles since
\((M+2)\ell<d\) eventually. If a sequence of \(M+1\) complete tiles
in the original indexing wraps, divide it where tile index zero is
reached. One of the two uninterrupted parts has at least
\(\lfloor M/2\rfloor\) tiles. On that part the parity of the original
coordinate is a fixed offset from the parity relative to its physical
start. The event list covers that start and that parity. Thus the
unlabelled remainder of length \(d\bmod\ell\), including the case of
odd \(d\), cannot destroy the marked tile.

It remains to intersect avoidance with the two global count conditions.
Let \(\mu=\sum_zp_z=pn(3o+1)/(4o)\). Each coset count has mean
\(\lambda_C\mu\). For independent Bernoulli sums of mean \(m\),
Bernstein's inequality gives failure probability at most
\(2\exp(-t^2/(2(m+t/3)))\) for absolute deviation exceeding \(t\).
With \(t=\mu x^{-1+\varepsilon}/(4(o+1))\), a union over the \(o\)
cosets gives failure probability at most
\(2o\exp(-c_o x^{79\varepsilon/8})\). On success
\(|n_T-\mu|\le ot\), which gives both bounds in \eqref{orb:cosets}.
Each overlap count has mean \(p_z n_2\) because \(T_2\subset A_0\).
For \(\lambda=(pn_2)^{1/2}x^{4\varepsilon}\) one has
\[
 pn_2\gg x^{2+55\varepsilon/8},
 \qquad \lambda/(pn_2)\ll x^{-1+9\varepsilon/16}=o(1).
\]
Bernstein and a union over the \(n\) translates give failure probability
at most \(2n\exp(-c' x^{8\varepsilon})\) for \eqref{orb:overlap-count}.
The sum of these two failure probabilities is smaller than the
displayed lower bound for avoidance. Their intersection with avoidance
therefore has positive probability.
\end{proof}

At the arithmetic host put \(h_q=\code{hPH}(q)\).
\(\code{Arith.Loc.Good}(B,h_q,a)\) applies precisely those three
conditions to the top positions from \(a\), with the original \(B.L_2\).
The prime set in \code{BadOK} is
\[
 \mathcal P=\{q<\lceil X'\rceil:\ q\ {\rm prime},\
       \code{IsResidueClassZero}(q),\
       \operatorname{ord}(h_q)\ge x^{\varepsilon/4}\}.
\]
With \(b_q(t)=\code{bsH}(\code{stepP}(B,h_q,q),t)\), its left side is
\[
 \sum_{q\in\mathcal P}\frac{\log q}{q}
   \#\{t\in T:\neg\code{Good}(B,h_q,b_q(t))\}.
\]
It equals zero for the set in Proposition~\ref{orb:selection}, even
when the specified origin depends on the selected data. The right side
is \(4C_b x^{-2}pn\sum_{q\in\mathcal P}\log q/q\), so any \(C_b\ge0\)
works. This verifies the \code{BadOK} clause of \code{HostGood}. On
a realization of this set, the original prime classification proves
that a missing top at a served prime must belong to this bad set.
Consequently \(T\setminus T_q=\varnothing\) at every served prime,
and its omitted-top charge vanishes. The first- and second-kind prime
losses, the net local lattice gains, and the realization of the
selected tops by the descended columns are separate hypotheses.

\begin{corollary}[Omitted tops in a previously chosen family]\label{orb:fixed-charge}
Let a family with its actual geometric records satisfy \code{ChosenE}.
There are constants \(C_b>0\) and \(x_0\) such that, for every
\(x\ge x_0\) and every index of that family at \(x\), its served-prime
charge satisfies
\[
 \sum_{q\ {\rm served}}4ma\frac{\log q}{q}\,|T\setminus T_q|
    \le 64 C_b(n_Tmb)x^{-1+\varepsilon/3}.
\]
The bound holds for every choice of the two auxiliary real parameters
\(C,b_o\) in the prime ledger.
\end{corollary}
\begin{proof}
Each missing top contributes to the bad count at the same prime by
the served-prime classification above. The total weighted count is
at most \(4C_b x^{-2}pn\sum_{q\in\mathcal P}\log q/q\) by the
original bad-top estimate. Here \(pn\le 2n_T\),
\(a=xb\) and \(\sum_{q\in\mathcal P}\log q/q\le 2\log X'
=2x^{\varepsilon/3}\). Multiplying by \(4ma\) gives the result.
The eventual top-count and bad-count conditions in \code{ChosenE}
use thresholds common to every index at the chosen scale.
\end{proof}

\paragraph{Second-kind primes below the local cutoff.}
Fix a candidate prime \(q<p_0(k)=k^2x^{26}\). For each good representative \(v\)
let \(N_v\) count its tops. These fibres are disjoint, so
\(\sum_vN_v\le n_T\). The small-prime definition gives, for every
choice of the two weight vectors,
\[
 0\le \operatorname{Sig}(q)
   =m\sum_v\min\{1,K_{\rm II}N_v r^2 k x^{-26}\}
   \le K_{\rm II}r^2 m n_T k x^{-26}.
\]
The original family bounds give \(m\le\sqrt x\),
\(n_T\le n\le x^{2+(8+1/100)\varepsilon}\) and
\(k\le x^{9/2+(10+1/100)\varepsilon}\). Thus
\(\operatorname{Sig}(q)\le K_{\rm II}r^2
x^{-19+(18+2/100)\varepsilon}\le x^{-18}\) past a single threshold
depending on the family and \(\varepsilon\). A second-kind prime in
the ledger must be a candidate with \(\operatorname{Sig}(q)>x^{-\varepsilon}\),
so all second-kind primes lie at or above \(p_0(k)\).
The residual share sum for candidates below the cutoff satisfies
\(\sum_q(\log q/q)4\operatorname{Sig}(q)\ll x^{-18}\log x\), since
\(p_0(k)\le x^{36}\). At \(q\ge p_0(k)\) the definition instead
uses the local quantity \code{kapI}, which depends on both weights.

\paragraph{The rank loss at a large prime.}
Fix a top and its pair of weights. Choose one admissible local ring
setting and one target lattice. Write \(p\) for the prime, \(w\)
for the uniformizer of the coefficient DVR \(\mathcal O\), with
\((p)=(w^2)\), and \(D\) for the prescribed depth. There are finitely
many dominant lifts \(\lambda\): dominance bounds the coordinates
between the prescribed endpoints on a top, and the other coordinates
have the prescribed cutoff. For each lift retain all actual source
columns satisfying the edge and off-edge conditions, and all vectors
in the prescribed filtered target lattice. Their pairings are integral.
At dimension \(d_\lambda\), a rank \(N\) certificate is equivalent to
an integral row combination of the evaluation matrix on finitely many
such columns and \(N\) target vectors satisfying
\[
        C\big(\langle F_i,y_j\rangle\big)_{ij}
                  \equiv p^D I_N\pmod{w^{2D+1}} .
\]
After division by \(p^D\) the error is in \(wM_N(\mathcal O)\)
and the resulting matrix is invertible over \(\mathcal O\).
The converse follows by evaluating the original interpolation on
the standard basis. The original ring and target transport lets one
use the same ring and lattice for all the lifts.
\src{src/Formalism/SharpRestriction/GrowthBounds/ColumnSpanDensity/RhoLowerBoundLargeConductor.lean}{138}{the local rank-loss construction}.
The companion \code{LargePrimePrecision.lean} checks the equivalence
with membership in the original \code{kapISet}.

On each side of the pairing take the integral span and divide by
its annihilator. The two images are finite free: independent actual
vectors on the opposite side give an injection into a finite power
of \(\mathcal O\), by integrality. Their common generic rank is the
rank of the pairing. A basis of the source image uses only finitely
many original columns, which can all be retained as generators.
Let \(s_{\lambda,i}\) be the Smith orders of this integral pairing,
padded by infinity to \(d_\lambda\). The maximal certificate rank is
\(\#\{i:s_{\lambda,i}\le 2D\}\). A shallow diagonal entry can be
scaled to \(p^D\). Conversely, if fewer than \(N\) entries are
shallow, a primitive integral row annihilates their columns in the
proposed \(N\)-row test. The remaining entries vanish modulo
\(w^{2D+1}\), whereas \(p^D\) times that row does not. Thus,
with \(l_\lambda(t)=\sum_i\min(t,s_{\lambda,i})\), the actual
infimum is attained and equals
\[
  \max\!\left(\{0\}\cup
       \left\{\frac{l_\lambda(2D+1)-l_\lambda(2D)}{d_\lambda}
                     : d_\lambda>0\right\}\right).
\]
The maximum is taken at each top and weight pair, before averaging
over the weight box. Estimating this maximum requires the filtered
pairings of the actual local columns.

\paragraph{A consecutive-window estimate for the Gram contribution.}
Let \(A(x)=\sum_{j=0}^{\delta}A_j\binom{x}{j}\) with
\(A_j\in M_N(\mathbf Z_p)\), \(N>0\), and \(P=\det A\ne0\) of
degree \(d_P\le N\delta\). For \(W\ge1\) consecutive integers \(z\),
let \(b_z\) count the Smith orders of \(A(z)\) at least \(K\ge1\),
including infinite orders. Normalize \(v(p)=1\) and put
\(h_\delta=\lfloor\log_p\max(1,\delta)\rfloor\). Then
\[
 \frac{1}{NW}\sum_z b_z
 \le \frac{\delta}{W}+
 \frac{\,\big(v(\operatorname{lc}P)+d_P/(p-1)\big)/N
                +\delta h_\delta/W\,}{K}.
\]
Indeed put \(m=K+h_\delta\). For \(1\le j\le\delta\) the identity
\[
  \binom{p^m X}{j}
     =(-1)^{j-1}\frac{p^m X}{j}
                    \prod_{i=1}^{j-1}(1-p^mX/i)
\]
and Vandermonde's formula give
\(A(z+p^mX)-A(z)\in p^K M_N(\mathbf Z_p[X])\).
Apply the integral Smith row and column transformations at \(z\).
In the perturbed matrix \(b_z\) rows are divisible by \(p^K\).
Thus the minimum coefficient valuation of \(P(z+p^mX)\) is at
least \(K b_z\), also when \(P(z)=0\). By multiplicativity of the
polynomial Gauss norm this minimum is
\(v(\operatorname{lc}P)+\sum_{\alpha}\min(m,v(z-\alpha))\),
with roots counted with multiplicity in a splitting field.
For \(j-1<t<j\), the integers with \(v(z-\alpha)>t\) lie in one
residue class modulo \(p^j\): the valuation of the difference of
two distinct integers is integral. They number at most \(W/p^j+1\).
Integrating for \(0<t<m\) bounds the mean contribution of each
root by \(1/(p-1)+m/W\), with negative valuations only lowering it.
Sum over the roots, divide by \(NK\), and use \(d_P\le N\delta\).

In the original Gram window \(A(z)\) is the recorded matrix of
the form multiplied by a nonzero integral scalar. This multiplication
can only increase the deep Smith count. The original leading-coefficient
and degree bounds now give its Gram estimate, whose normalized cost
at the original depths is at most \(5/Y\).
\src{src/Formalism/SharpRestriction/GrowthBounds/FockGramWindow/GramWindow2.lean}{12}{Gram window hypotheses}
and \src{src/Formalism/SharpRestriction/GrowthBounds/FockGramWindow/InitialPointNumericBound.lean}{235}{the normalized bound}.
This estimates the Gram contribution. The harmonic ranks, the
off-edge restrictions on the actual columns and the required
uniform weight-box bound on the maximum above still enter the
large-prime loss.

\begin{corollary}[Size of the overlap error]
Use the common Fourier weights\label{orb:overlap} and thin-coordinate estimate in the
translate comparison. Put \(B=n_Tmb\) and \(\delta_2=n_2/n\).
Let \(w=\max_{t\in T}|a_t/b-x|\le1\), and suppose the kernels satisfy
\[
   \sum_h|K_h(x)|\ll xL,\qquad
   \sum_h\sup_{|z-x|\le w}|K_h(z)-K_h(x)|\ll wL,
   \qquad L\ll x^{\varepsilon/3}.
\]
Writing \(S_{\psi;q,d}(x)\) for the signed four-twist sum, the same
set \(T\) then satisfies
\[
  \sup_{u\in A_0}\left|\frac{2D_u}{B\delta_2}-S_{\psi;q,d}(x)\right|
   \ll x^{43\varepsilon/48}+w x^{\varepsilon/3}
                   + (x\delta_2)^{-1}
   =o(\sqrt x).
\]
\end{corollary}
\begin{proof}
Thick replacement and thin deletion cost at most \(2B/x\) in \(D_u\).
In the remaining convolution the multiplicity at offset \(h\) is
\(|T\cap(u-h+T_2)|\). Moving the endpoints to \(x\) costs at most
\(O(mb(pn_2+\lambda)wL)\). The resulting discrepancy from the
convolution with multiplicity \(n_2p_{-h}\) is \(O(mb\lambda xL)\)
by \eqref{orb:overlap-count}. Summing the overlap errors over \(z\)
also gives \(|n_T-\mu|\le n\lambda/n_2\). The signed four-twist sum
is \((2n/\mu)\sum_hp_{-h}K_h(x)=O(xL)\) by finite character
orthogonality. After normalizing, the total cost is
\(O(xL\lambda/(pn_2)+wL+1/(x\delta_2))\). Substitution gives
\(1-1=0\) in the main power of \(x\) and
\((1/3+9/16)\varepsilon=43\varepsilon/48\) in the remaining power.
\end{proof}

\begin{corollary}[Overlap at the binomial scale]\label{orb:conditional-overlap}
Condition the Bernoulli selection on avoiding all the local events of
Proposition~\ref{orb:selection}. With conditional probability
\(1-o(1)\), its coset and cardinality bounds hold and, simultaneously,
\[
 \max_{z\in H}\big||T\cap(z+T_2)|-p_z n_2\big|
    \ll \sqrt{pn_2\log(2n)}.
\]
The overlap contribution in Corollary~\ref{orb:overlap} is then
\(O(x^{-149\varepsilon/48}\sqrt{\log(2x)})\). The endpoint and
thin-coordinate terms remain \(O(wL+1/(x\delta_2))\).
\end{corollary}
\begin{proof}
Write \(\mathcal G\) for avoidance and put
\(\eta=(1-2q_*)^{-4nG'}-1=O(x^{-2+\gamma/2}\log x)\).
A cylinder event \(A\) specifying the states of \(j\) points meets
the supports of at most \(4jnG'\) local events. It is independent of
joint avoidance of all the others. Conditioning first on those others,
the chain rule and the conditional bound \(2q_*\) from the proof of
Proposition~\ref{orb:selection} give
\(\Pr(A\mid\mathcal G)\le(1+\eta)^j\Pr(A)\). This holds for
every assignment, including all zeros and the empty assignment.

Let \(S\) be one coset or one translate of \(T_2\), with \(s=|S|\)
and common inclusion probability \(\theta\in[3p/4,p]\).
Set \(Y=|T\cap S|\). For \(u\ge0\), expand \(e^{uY}\) as the product
of \(1+(e^u-1)\mathbf1_{z\in T}\) over \(z\in S\). All coefficients
are nonnegative. The cylinder bound applied to each monomial gives
\[
 \mathbb E(e^{uY}\mid\mathcal G)
 \le [1+(e^u-1)(1+\eta)\theta]^s.
\]
The same expansion for the unselected points gives this inequality
for \(s-Y\) with \(\theta\) replaced by \(1-\theta\). Since
\(\eta=o(p)\), both right sides are binomial moment generating
functions with variance at most \(2ps\). Each centre moves by at most
\(\eta s\). Uniformly for \(0<s\le n\) we have
\(\eta\sqrt{s/p}\ll x^{-1+4\varepsilon+\gamma}\log x=o(1)\).
For the sets here \(ps\gg x^{2+55\varepsilon/8}\). The exponential
Markov bound for each of the two binomial moment generating functions
thus bounds the probability of deviation from \(s\theta\) exceeding
\(C\sqrt{ps\log(2n)}\) by \(O(n^{-3})\) for a sufficiently large
absolute \(C\). Sum this over the \(n\) translates and the \(o\)
cosets. For a coset this deviation is \(o(pn x^{-1+\varepsilon})\),
so the argument for \eqref{orb:cosets} applies. Finally
\(xL\sqrt{\log(2n)/(pn_2)}\ll
x^{-149\varepsilon/48}\sqrt{\log(2x)}\) by the same lower bound on
\(pn_2\).
\end{proof}

The original finite conditions appear in \code{GoodOrbit\allowbreak
IntervalTile\allowbreak Condition.lean}, lines 10--26, with the cyclic
window in \code{StatementB.lean}, lines 9--13. The parameters appear
in \code{Unramified\allowbreak Quadratic\allowbreak Setting.lean}, lines
239--252. The arithmetic host uses these conditions in
\code{GlobalTheory/\allowbreak IsResidueClassZero.lean}, lines 96--104.
The prime sum appears in \code{ThinScales\allowbreak IffHostBound.lean},
lines 36--42 and 101--104. The served-prime classification is proved
in \code{PowFourLt\allowbreak CeilThreshold\allowbreak Prime.lean},
lines 722--761. The companion \code{MeanLoss\allowbreak Refinement.lean}
checks Corollary~\ref{orb:fixed-charge} for the original family types.
The small-prime share and its cutoff occur in
\code{GrowthBounds/\allowbreak IsResidueClassZero.lean}, lines 202--203,
227--228, 276--277 and 287--288. The orbit counts are identified and summed in
\code{SlopeEstimateOf\allowbreak LemmaBSlopeLine.lean}, lines 129--146.
